% CR10000C
% Akutes Abdomen
% 14.0
% 2013-11-30

\label{m:akutes-abdomen}
\EE{Taktik}: \Speech{Unmittelbare vitale Bedrohung
einschätzen, sie ist
abhängig von den Vitalwerten.
Vermutliche Ursache und Differentialdiagnosen
abklären, Transport an eine geeignete Abteilung}


\begin{itemize}
  \item Allgemeine Maßnahmen bei Abdominalerkrankungen
  (\prettyref{m:am-abdominalerkrankungen})
  \item \EE{Vitale Bedrohung einschätzen.} Ggfs.
  \TxMassVitMK
  \item Lagerung: grundsätzlich \EE{bauchdeckenentspannend}
  (Knierolle, Beine angezogen), oder nach Wunsch des Patienten
  \item \EE{Nüchtern lassen} (in Hinblick auf möglicherweise
  bevorstehende Operation, nichts essen oder trinken lassen)
  \item Monitoring: Der Zustand kann sich abrupt
  verschlechtern.
  \item \EE{Transportentscheidung}: Abt. f. \EE{Chirurgie},
  bei Kindern ${\rightarrow}$ Abt. f. \EE{Kinderheilkunde},
  bei Frauen evtl. \EE{Gynäkologie} oder Chirurgie mit Abt.
  f. Gynäkologie im Hintergrund
\end{itemize}